% Part: lambda-calculus
% Chapter: syntax
% Section: de-bruijn

\documentclass[../../../include/open-logic-section]{subfiles}

\begin{document}

\olfileid{lam}{syn}{deb}

\olsection{डी ब्रुइन निर्देशांक}

$\alpha$-सममूल्यता ही स्वाभाविक कल्पना आहे, कारण
$\alpha$-सममूल्य पदांचा ``अर्थ एकच'' असतो. प्रत्यक्षात,
एकमेकांत $\alpha$-परिवर्तन करता येणारी पदे वेगळी
न मानणारा लॅम्डा पदांचा विन्यास देता येतो. यांपैकी
सर्वाधिक परिचित पद्धतीत चरांच्या जागी त्यांचे
\emph{डी ब्रुइन निर्देशांक} वापरतात.

$\lambd[x][M]$ लिहिताना फलनाचा प्राचल $x$ आहे हे
स्पष्ट सांगतो, म्हणजे~$M$ मध्ये त्या प्राचलाचा संदर्भ
घेण्यासाठी $x$ वापरता येतो. मात्र डी ब्रुइन निर्देशांक
पद्धतीत प्राचलांना नावे नसतात. फलनाच्या मुख्य भागात
प्राचलाचा संदर्भ हा त्याच्या आस्थितीपासून त्याला बद्ध
करणाऱ्या अमूर्तीकरणापर्यंत मधे येणाऱ्या अमूर्तीकरणांची
पातळी मोजणाऱ्या संख्येने दर्शवतात. उदाहरणार्थ,
$\lambd[x][\lambd[y][y x]]$ पाहा: बाहेरील अमूर्तीकरण
चर~$x$ बद्ध करते; आतील अमूर्तीकरण चर~$y$ बद्ध करते.
उपपद $y x$ हे आतील अमूर्तीकरणाच्या व्याप्तीत आहे.
$y$ आणि त्याला बद्ध करणारे~$\lambd[y]$ यांच्यात
एकही अमूर्तीकरण नाही; पण $x$ आणि त्याला बद्ध
करणारे~$\lambd[x]$ यांच्यात एक अमूर्तीकरण आहे.
म्हणून $y x$ साठी $0\, 1$, आणि संपूर्ण पदासाठी
$\lambd[][\lambd[][01]]$ असे लिहितात.

\begin{defn}
  डी ब्रुइन पदे आगमनाने पुढीलप्रमाणे परिभाषित करतात:
  \begin{enumerate}
  \item $n$, जेथे $n$ ही कोणतीही नैसर्गिक संख्या आहे.
  \item $PQ$, जेथे $P$ आणि $Q$ ही दोन्ही डी ब्रुइन पदे आहेत.
  \item $\lambd[][N]$, जेथे $N$ हे डी ब्रुइन पद आहे.
  \end{enumerate}
\end{defn}

सामान्य लॅम्डा पदांचे डी ब्रुइन निर्देशांकित पदांत
औपचारिक रूपांतर पुढीलप्रमाणे करतात:
\begin{defn}
  \begin{align*}
    F_\Gamma(x) &= \Gamma(x) \\
    F_\Gamma(PQ) &= F_\Gamma(P)F_\Gamma(Q) \\
    F_\Gamma(\lambd[x][N]) &= \lambd[][F_{x,\Gamma}(N)]
  \end{align*}
  येथे $\Gamma$ ही शून्यापासून निर्देशांकित चरांची यादी
  आहे आणि $\Gamma(x)$ ही $\Gamma$ मधील चर $x$ ची
  जागा दर्शवते. उदाहरणार्थ, $\Gamma$ ही $x,y,z$
  असेल, तर $\Gamma(x)$ हे $0$ आणि $\Gamma(z)$ हे $2$ आहे.

  $x,\Gamma$ म्हणजे $\Gamma$ च्या सुरुवातीला $x$
  ठेवून मिळालेली यादी. उदाहरणार्थ, मागील उदाहरणातील
  $\Gamma$ साठी $w,\Gamma$ ही $w,x,y,z$ आहे.
\end{defn}

डी ब्रुइन पदापासून नेहमीचे लॅम्डा पद पुढीलप्रमाणे
परत मिळवतात:

\begin{defn}
  \begin{align*}
    G_\Gamma(n) &= \Gamma[n] \\
    G_\Gamma(PQ) &= G_\Gamma(P) G_\Gamma(Q) \\
    G_\Gamma(\lambd[][N]) &= \lambd[x][G_{x,\Gamma}(N)]
  \end{align*}
  येथेही $\Gamma$ ही शून्यापासून निर्देशांकित चरांची
  यादी आहे, आणि $\Gamma[n]$ म्हणजे~$n$ व्या जागेवरील
  चर. उदाहरणार्थ, $\Gamma$ ही $x,y,z$ असेल, तर
  $\Gamma[1]$ हे $y$ आहे.

  शेवटच्या समीकरणातील चर $x$ हा $\Gamma$ मध्ये
  नसलेला कोणताही चर निवडतात.
\end{defn}

येथे काही निष्कर्ष त्यांच्या सिद्धता न देता मांडतो:

\begin{prop}
  $M \aconvone M'$ असेल, आणि $\Gamma$ ही
  $\FV{M}$ समाविष्ट करणारी कोणतीही यादी असेल,
  तर $F_\Gamma(M) \eqs F_\Gamma(M')$.
\end{prop}

\end{document}
